\documentclass[12pt,a4paper]{article}
\usepackage[utf8]{inputenc}
\usepackage{amsmath, amssymb, amsfonts}
\usepackage{geometry}
\geometry{margin=2.5cm}
\usepackage{hyperref}
\usepackage{pgfplots}
\usepackage{fancyhdr}
\usepackage{listing}
\usepackage{listings}
\usepackage{titling}
\usepackage{float}
\usepackage{titlesec}

\titleformat{\section}
  {\normalfont\Large\bfseries\centering}  % Format: font size, bold, centered
  {\thesection}{1em}{}
\pagestyle{fancy}
\pgfplotsset{compat=1.18}
\fancyhead{}
\fancyfoot{}
\title{Il colloquio di lavoro}
\pretitle{\begin{center}\LARGE}
\posttitle{\par\large Career Labs\end{center}}
\author{}
\date{}
\fancyhead[L]{Gabriele Perna}          % Left side
\fancyhead[C]{Il colloquio di lavoro}     % Center
\fancyhead[R]{2026/2027}             % Right side
\fancyfoot[C]{\thepage}           % Page number in the center of footer
\let\oldsection\section
\renewcommand{\section}{\clearpage\oldsection}

\begin{document}

\maketitle
\tableofcontents

\section{Informazioni e risorse}
\textbf{Nominativo consulente:} Giulia Verdigi.\\\\
\textbf{Consulenza} \textbf{online} \textbf{(gratuita) sui colloqui di lavoro:}
\href{https://alumni.unipi.it/mentoring}{https://alumni.unipi.it/mentoring}

\section{Prepararazione al colloquio}
Le cose principali di cui occuparsi nel periodo prima del colloquio sono:
\begin{itemize}
    \item Informazioni sull'azienda (sito, social...)
    \item Orario del colloquio
    \item Controllare il luogo di lavoro
    \item Preparazione del CV
    \item Look e outfit
    \item Preparazione psico-fisica
    \item Comprendere la posizione offerta nel dettaglio
    \item Capire l'organigramma
    \item Capire e conoscere i competitor
    \item Preparazione di domande da fare al recruiter
    \item Preparazione dello \textbf{speech}
    \item Il \textbf{CNL} (contratto nazionale di lavoro) per confrontarlo con la retribuzione offerta
\end{itemize}
\section{Fare bella impressione}
\subsection{La chiamata telefonica}
In particolare, la prima fase di selezione di un colloquio non avviene durante il colloquio vero e proprio, bensì durante la prima \textbf{chiamata telefonica} o il primo scambio di messaggi \textbf{email}.\\
Durante la telefonata bisogna essere svegli e prendere appunti, non bisogna perdere informazioni importanti.
\subsection{Puntualità e attesa}
Al colloquio di lavoro bisogna arrivare puntuali e bisogna prevedere ogni imprevisto, è consigliato svegliarsi e partire assai prima dell'orario necessario.\\
L'orario di arrivo non dipende esclusivamente da te, però un ritardo o un'assenza ricade direttamente su di te lo stesso.\\
L'orario giusto a cui presentarsi può variare da vari fattori, come la presenza di una sala d'attesa o l'apertura della struttura.\\\\
Oltre a questo, non si deve assillare il recruiter con messaggi che specificano che siete arrivati e bisogna annoiarsi, non si deve usare troppo il telefono.
\subsection{Dress e web code}
Il modo in cui ci si presenta dice molto su di noi, non solo nel modo in cui ci si esprime coi vestiti, ma anche per esempio con il nostro igiene.\\
Per fare bella figura, bisogna essere puliti e presentarsi bene.\\
Abbiamo il diritto di vestirci come ci pare, bisogna però avere una strategia in mente che agisca a nostro favore.
\subsection{Accompagnamenti}
Bisogna evitare assolutamente di essere accompagnati dai genitori, partner o animali domestici, bisogna presentarsi da soli per dimostrare la nostra indipendenza.
\subsection{Linguaggio fisico}
Essere aperti aiuta nella conversazione, spalle dritte, composti sulla sedia.\\
Stringere la mano al recruiter e parlare ad alta voce.\\

\section{Domande vietate}
Ecco alcune delle domande che un recruiter non può farvi \textbf{legalmente}:
\begin{itemize}
    \item Qual'è il tuo stato di salute?
    \item Sei disabile?
    \item Hai mai sofferto di disabilità?
    \item A quale religione appartieni?
    \item Quali festivita segui?
    \item Da quale nazione provieni?
    \item Sei sposato/a?
    \item Sei incinta?
    \item Stai programmando una gravidanza?
    \item Hai intenzione di avere figli?
    \item Hai figli? Se si, come concilieresti gli impegni?
    \item Sei mai stato arrestato?
    \item Hai debiti?
    \item Fai uso di sostanze stupefacenti?
    \item Bevi alcolici?
    \item Quanto guadagnavi prima?
\end{itemize}
A queste domande possiamo non rispondere o semplicemente evitare di dire la verità.

\section{Consigli importanti}
Uno dei consigli più importanti per un colloquio è di \textbf{non mentire}, bisogna essere se stessi e non svalutare le nostre esperienze.\\
Allo stesso tempo, è consigliato valorizzare gli hobby senza renderli tema centrale del colloquio.

\end{document}